\section{Оценка эффективности предложенного подхода}

Разработанный гибридный подход позволяет решить задачу восстановления истории построения с приемлемой точностью для основных операций. Полученные результаты (таблица~\ref{tab:seg-results}) демонстрируют, что:
\begin{itemize}
    \item для операций \texttt{fillet} и \texttt{cutExtrusion} точность сегментации превышает 90\%;
    \item классификатор для \texttt{fillet} достигает 90.6\% accuracy;
    \item даже для менее представленной операции \texttt{chamfer} точность выше случайного угадывания.
\end{itemize}

Преимущества предложенного решения:
\begin{itemize}
    \item Модульность: можно обучать и добавлять новые операции независимо.
    \item Интерпретируемость: сегментация позволяет визуально контролировать результат.
    \item Интеграция с CAD-ядром: точные параметры вычисляются детерминированно, что повышает надёжность.
\end{itemize}

Ограничения:
\begin{itemize}
    \item Зависимость от качества датасета (дисбаланс, сложность моделей).
    \item STL теряет инженерную семантику, что ограничивает точность параметрической реконструкции.
    \item В текущей версии поддерживаются только четыре типа операций.
\end{itemize}

Тем не менее, полученные результаты подтверждают принципиальную осуществимость подхода и его пригодность для дальнейшего масштабирования.

